\section{ACT 近似算法伪代码}

\begin{algorithm}[t]
\caption{ACT 近似。}
\begin{footnotesize}
\begin{algorithmic}[1]
\State 保留原文伪代码结构，用于说明近似目标函数的计算流程
\State 其核心思想是：精确计算本轮候选动作的 ACT，并对后续动作的 ACT 进行近似估计
\end{algorithmic}
\end{footnotesize}
\label{app:act_approx}
\end{algorithm}

\section{统一 DPArrange 算法}

本节给出 DPArrange 的附录说明。该算法负责在给定资源约束下，为候选动作求解离散资源分配，并可适配不同拓扑结构。

\begin{algorithm}[t]
\caption{与拓扑无关的 DPArrange}
\begin{footnotesize}
\begin{algorithmic}[1]
\State 保留原文算法主体所表达的动态规划思想
\State 目标是在候选动作之间寻找满足离散资源约束的最优资源划分
\end{algorithmic}
\end{footnotesize}
\label{app:udp}
\end{algorithm}

\begin{algorithm}[t]
\caption{GPU 拓扑感知 DP 算子}
\begin{footnotesize}
\begin{algorithmic}[1]
\State 该算子在统一 DPArrange 之上，进一步处理 GPU 拓扑相关约束
\State 包括并行度、设备组选择与资源布局等因素
\end{algorithmic}
\end{footnotesize}
\label{app:gpu_dp}
\end{algorithm}
